\chapter*{摘要}
\addcontentsline{toc}{chapter}{摘要}

低功耗广域物联网中，LoRa 终端的大规模部署使设备身份认证成为网络安全的关键环节。射频指纹识别（RFFI）利用发射机硬件非理想性实现物理层身份验证，可作为密码认证的补充。然而，真实部署中存在跨日期信道漂移、跨接收机前端失配、复杂电磁扰动以及未知设备入网等问题，使 LoRa RFFI 的识别与认证可靠性面临严峻挑战。

本文围绕上述挑战，从建模、跨域校准与复杂环境评估三个层次开展研究。（1）提出带外（OOB）频谱引导的 cross-attentive RF-HSTU 混合模型，融合带内 IQ 与带外谱证据，在同接收机跨日期闭集协议下文件级准确率达 $75.0\pm5.3\%$，较 CNN-IQ 基线 $54.2\pm14.2\%$ 提升约 20.8 个百分点。（2）诊断跨接收机迁移失败机理，揭示 OOB 路径上的接收机诱导特征纠缠；提出 RCPA-T 目标接收机原型校准方法，在 block-disjoint 少样本协议下 $K{=}20$ 时达 $75.0\%$，显著优于 source-only 近随机水平。（3）构建复杂电磁扰动 benchmark，扩展开放集未知设备认证协议；证明 Prototype/Mahalanobis 嵌入空间评分（AUROC $0.917\pm0.059$）优于 MSP/Energy；发现载波频偏（CFO）同时破坏已知类分类与未知拒识，为当前系统最强失效模式。扰动一致性训练（EM-CR）初步探索未达主方法标准，作为未来工作讨论。

本文工作为 LoRa RFFI 在复杂部署条件下提供了可复现的评估协议、诊断框架与性能边界分析，而非声称已完全解决设备认证问题。

\vspace{1em}
\noindent\textbf{关键词：} LoRa；射频指纹识别；设备认证；带外频谱；跨接收机校准；开放集认证；电磁鲁棒性
